\chapter{Newton's laws and applications}\label{ch:newton}

Kinematics (chapter~1) described motion; \emph{dynamics} explains it.
Newton's three laws are exactly three statements --- a definition of
inertial frames, a relation $\mathbf F=m\mathbf a$, and a pairing rule ---
but the engineering content is in the method that wraps them: the free-body
diagram, the constraint bookkeeping of strings and surfaces, and the
differential equations that friction and drag leave behind.

% from-zero:begin
\begin{rulebox}[Before you start]
Chapter~1: position, velocity, acceleration. \emph{Recap:} acceleration is the rate of change of velocity. Newton's laws say what causes acceleration.
\end{rulebox}
\begin{wins}[After this chapter you can]
\begin{itemize}
\item state the three laws and the frame in which they hold;
\item draw a free-body diagram for any situation;
\item solve inclines, pulleys and circular motion;
\item handle friction and drag.
\end{itemize}
\end{wins}
\begin{principle}[How to think about this]
Newton's second law is a recipe: find every force on one object, add them as vectors, and divide by the mass. The free-body diagram is how you keep the list honest. If you draw it every time, even hard problems become the same routine.
\end{principle}
% from-zero:end

\section{The three laws, and the frame assumption}\label{sec:new-laws}

\textbf{First law.} A body with no net force moves at constant velocity ---
\emph{in an inertial frame}. The law looks circular (how do you know the
frame is inertial? check whether first-law behaviour holds) and it is:
Newton's first law is the \emph{definition} of the arena his other two laws
live in. Accelerating cars and rotating merry-go-rounds are non-inertial;
inside them you feel ``forces'' (centrifugal, elevator heaviness) that have
no third-law partner.

\textbf{Second law.}
\begin{equation}
  \sum\mathbf F=m\mathbf a.\label{eq:fma}
\end{equation}
This is a \emph{vector} equation --- three scalar equations when you pick
axes --- and the most compressed sentence in physics: acceleration is what
the \emph{sum} of forces buys. Units: force in newtons,
$1\,\mathrm{N}=1\,\mathrm{kg\cdot m/s^{2}}$.

\textbf{Third law.} Forces come in pairs on \emph{different} bodies: the
Earth pulls you down; you pull the Earth up. The pair never cancels on one
body's diagram --- a point section~\ref{sec:new-fbd} returns to.

\section{The force inventory}\label{sec:new-inventory}

Only a handful of forces appear in this book:

\begin{center}
\begin{tabularx}{\linewidth}{@{}llY@{}}
\toprule
\tablehead\thead{Force} & \thead{Rule} & \thead{Notes}\\
\midrule
gravity & $m\mathbf g$ down & uniform near Earth's surface \\
normal & whatever it must be & contact; perpendicular to surface \\
tension & along the rope & pulls both ends; massless rope $\Rightarrow$ same $T$ \\
spring & $-kx$ & Hooke's law, linear \\
friction & $\le\mu_sN$ static, $=\mu_kN$ kinetic & opposes \emph{impending/actual} sliding \\
drag & $cv$ or $cv^{2}$ & opposes velocity; \S\ref{sec:new-drag} \\
\bottomrule
\end{tabularx}
\end{center}

The discipline: every entry except gravity is a \emph{contact} force ---
if nothing touches the body, only $m\mathbf g$ acts. ``The force of the
throw'' left the ball when your hand did.

\section{Free-body diagrams: the discipline}\label{sec:new-fbd}

A free-body diagram (FBD)\index{free-body diagram} draws \emph{one} body and
every force \emph{on} it --- nothing else. Not the force it exerts, not its
velocity, not its acceleration (draw $m\mathbf a$ beside the diagram, never
in it). Figure~\ref{fig:fbdgal} is the canonical gallery.

\begin{figure}[htbp]\centering
\begin{tikzpicture}[scale=1.0,force/.style={->,thick,accent}]
% 1: free fall
\begin{scope}[xshift=0cm]
  \filldraw[fill=gray!20] (0,0) rectangle (0.5,0.5);
  \draw[force] (0.25,0.25) -- (0.25,-0.8) node[below,font=\scriptsize]{$m\mathbf g$};
  \node[font=\scriptsize] at (0.25,-1.2) {free fall};
\end{scope}
% 2: on table
\begin{scope}[xshift=1.9cm]
  \draw[thick] (-0.3,-0.15) -- (0.8,-0.15);
  \filldraw[fill=gray!20] (0,0) rectangle (0.5,0.5);
  \draw[force] (0.25,0.25) -- (0.25,-0.8) node[below,font=\scriptsize]{$m\mathbf g$};
  \draw[force] (0.25,0.5) -- (0.25,1.3) node[above,font=\scriptsize]{$\mathbf N$};
  \node[font=\scriptsize] at (0.25,-1.2) {table};
\end{scope}
% 3: hanging mass
\begin{scope}[xshift=3.8cm]
  \draw[thick] (-0.2,1.6) -- (0.7,1.6);
  \draw (0.25,1.6) -- (0.25,0.5);
  \filldraw[fill=gray!20] (0,0) rectangle (0.5,0.5);
  \draw[force] (0.25,0.25) -- (0.25,-0.8) node[below,font=\scriptsize]{$m\mathbf g$};
  \draw[force] (0.25,0.5) -- (0.25,1.3) node[above,font=\scriptsize]{$\mathbf T$};
  \node[font=\scriptsize] at (0.25,-1.2) {hanging};
\end{scope}
% 4: pushed with friction
\begin{scope}[xshift=5.7cm]
  \draw[thick] (-0.3,-0.15) -- (0.8,-0.15);
  \filldraw[fill=gray!20] (0,0) rectangle (0.5,0.5);
  \draw[force] (0.25,0.25) -- (0.25,-0.8) node[below,font=\scriptsize]{$m\mathbf g$};
  \draw[force] (0.25,0.5) -- (0.25,1.3) node[above,font=\scriptsize]{$\mathbf N$};
  \draw[force] (-0.5,0.25) -- (0,0.25) node[midway,above,font=\scriptsize]{$\mathbf F$};
  \draw[force] (0,0.1) -- (-0.4,0.1) node[above left=-2pt,font=\scriptsize]{$\mathbf f$};
  \node[font=\scriptsize] at (0.25,-1.2) {pushed};
\end{scope}
% 5: spring
\begin{scope}[xshift=7.6cm]
  \draw[thick] (-0.3,-0.15) -- (0.9,-0.15);
  \filldraw[fill=gray!20] (0.2,0) rectangle (0.7,0.5);
  \draw[thick] (-0.2,0.25) -- (-0.13,0.35) -- (-0.06,0.15) -- (0.01,0.35) -- (0.08,0.15) -- (0.15,0.35) -- (0.2,0.25);
  \draw[force] (0.45,0.25) -- (0.45,-0.8) node[below,font=\scriptsize]{$m\mathbf g$};
  \draw[force] (0.45,0.5) -- (0.45,1.3) node[above,font=\scriptsize]{$\mathbf N$};
  \draw[force] (0.7,0.4) -- (1.3,0.4) node[right,font=\scriptsize]{$k\mathbf x$};
  \node[font=\scriptsize] at (0.45,-1.2) {spring};
\end{scope}
\end{tikzpicture}

\caption{Five canonical FBDs. Notice only the body's own forces appear ---
the Earth's pull counts, the body's pull on the Earth does not. The pushed
box's friction arrow points \emph{back} because friction opposes sliding.}
\label{fig:fbdgal}
\end{figure}

The method: (1) isolate one body; (2) enumerate contacts + gravity; (3) pick
axes that simplify components; (4) write equation~\eqref{eq:fma} per axis;
(5) add constraint equations (rope length, contact, $f\le\mu_sN$); (6) solve.
Skipping step 2's enumeration is where forces get invented or dropped.

\begin{example}[Atwood machine]\label{ex:atwood}
Masses $m_1=3\,\mathrm{kg}$ and $m_2=5\,\mathrm{kg}$ hang from a light
frictionless pulley (figure~\ref{fig:atwood}). Find the acceleration and
tension.
\par\smallskip
\textbf{Strategy.} Two FBDs, one constraint: the rope's fixed length means
the accelerations are equal and opposite. Take each mass's ``down is
positive'' axis along its own motion.
\begin{solution}
For $m_1$ (accelerating up): $T-m_1g=m_1a$.
For $m_2$ (accelerating down): $m_2g-T=m_2a$.
Add to eliminate $T$:
\[
  a=\frac{m_2-m_1}{m_1+m_2}g=\frac{2}{8}\times9.8=2.45\,\mathrm{m/s^{2}},
  \qquad
  T=m_1(g+a)=3\times12.25=36.75\,\mathrm{N}.
\]
Check limits: $m_1=m_2\Rightarrow a=0,T=mg$ (hanging balance);
$m_2\gg m_1\Rightarrow a\to g$ (small mass is a brake). $T=36.75$\,N lies
between $m_1g=29.4$ and $m_2g=49$ --- the moving rope always splits the
difference. \checkmark
\end{solution}
\end{example}

\begin{figure}[htbp]\centering
\begin{tikzpicture}[scale=1.0,force/.style={->,thick,accent}]
\draw[thick] (-1.5,3.2) -- (1.5,3.2);
\draw[fill=white] (0,3.2) circle (0.35);
\draw (0,3.55) -- (0,3.2); \fill (0,3.2) circle (0.03);
\draw (-0.35,3.2) -- (-0.35,1.2);
\draw (0.35,3.2) -- (0.35,0.4);
\filldraw[fill=gray!20] (-0.7,0.7) rectangle (0,1.2);
\node[font=\scriptsize] at (-0.35,0.95) {$m_1$};
\filldraw[fill=gray!20] (0,0.0) rectangle (0.7,0.4);
\node[font=\scriptsize] at (0.35,0.2) {$m_2$};
\draw[force] (-0.35,0.7) -- (-0.35,-0.2) node[below,font=\scriptsize]{$m_1\mathbf g$};
\draw[force] (-0.35,1.2) -- (-0.35,2.1) node[above,font=\scriptsize]{$\mathbf T$};
\draw[force] (0.35,0.0) -- (0.35,-0.9) node[below,font=\scriptsize]{$m_2\mathbf g$};
\draw[force] (0.35,0.4) -- (0.35,1.3) node[above,font=\scriptsize]{$\mathbf T$};
\draw[->,thick,amber] (-0.95,1.6) -- (-0.95,2.4) node[above,font=\scriptsize]{$\mathbf a$};
\draw[->,thick,amber] (1.0,0.7) -- (1.0,-0.1) node[below,font=\scriptsize]{$\mathbf a$};
\node[font=\scriptsize] at (0,3.2) {};
\end{tikzpicture}

\caption{Atwood machine: two FBDs, one rope constraint. Notice the tension
is the \emph{same} $T$ on both sides --- the massless-frictionless ideal
pulley transmits, it does not eat.}
\label{fig:atwood}
\end{figure}

\begin{exercise}[easy]
In example~\ref{ex:atwood}, what is the force the pulley's axle must exert
on the ceiling? Why isn't it $m_1g+m_2g$? (AP-C)
\end{exercise}

\section{Constraints: inclines, pulleys, contact}\label{sec:new-incline}

On an incline the useful axes are \emph{down-slope} and \emph{outward}
(figure~\ref{fig:incline}): gravity decomposes as $mg\sin\theta$ down-slope,
$mg\cos\theta$ into the surface, and when the surface doesn't accelerate,
$\mathbf N=mg\cos\theta$ \emph{only then}. That ``only then'' is the
section's lesson --- $N$ is determined, not given.

\begin{figure}[htbp]\centering
\begin{tikzpicture}[scale=1.15,force/.style={->,thick,accent}]
\draw[thick] (0,0) -- (4.5,0) -- (4.5,2.2) -- cycle;
\draw (3.9,0) arc (180:180-26:0.6);
\node[font=\scriptsize] at (3.55,0.22) {$\theta$};
\begin{scope}[rotate around={-26:(2.3,0.66)}]
  \filldraw[fill=gray!20] (2.0,0.66) rectangle (2.6,1.16);
\end{scope}
\draw[force] (2.35,0.95) -- (2.35,-0.6) node[below,font=\scriptsize]{$m\mathbf g$};
\draw[force,dashed,amber] (2.35,0.95) -- (1.6,0.28) node[below left=-2pt,font=\scriptsize]{$mg\sin\theta$};
\draw[force,dashed,amber] (2.35,0.95) -- (2.95,1.65) node[above right=-2pt,font=\scriptsize]{$mg\cos\theta$};
\draw[force] (2.55,1.3) -- (3.2,2.0) node[above,font=\scriptsize]{$\mathbf N$};
\draw[force] (2.15,1.0) -- (1.45,1.35) node[above left=-2pt,font=\scriptsize]{$\mathbf f$};
\draw[->,dashed] (4.7,1.0) -- (5.7,1.0);
\begin{scope}[xshift=5.2cm,yshift=0.5cm]
  \draw[->] (0,0) -- (0.8,-0.37) node[right,font=\scriptsize]{down-slope $+$};
  \draw[->] (0,0) -- (0.34,0.74) node[above,font=\scriptsize]{outward};
\end{scope}
\end{tikzpicture}

\caption{Incline decomposition. Notice the axes are rotated, not the force:
$mg$ splits into a down-slope driver $mg\sin\theta$ and a pressed-in part
$mg\cos\theta$ that the normal force answers.}
\label{fig:incline}
\end{figure}

\begin{example}[Apparent weight on an accelerating surface]\label{ex:apparent}
A $60\,\mathrm{kg}$ student stands on a scale in an elevator accelerating
up at $2\,\mathrm{m/s^{2}}$. What does the scale read? What if the cable
snaps?
\par\smallskip
\textbf{Strategy.} The scale reads the \emph{normal} force, not gravity.
One FBD, one axis (up $+$).
\begin{solution}
$N-mg=ma \Rightarrow N=m(g+a)=60\times11.8=708\,\mathrm{N}$ --- the scale
reads 708\,N ``of weight'', i.e.\ a $72.3\,\mathrm{kg}$ equivalent.
Free fall $a=-g$: $N=0$, weightlessness --- not ``no gravity'', just
nothing to push against. The ISS floats exactly this way, inside $0.9g$ of
gravity. \checkmark
\end{solution}
\end{example}

\begin{exercise}[medium]
A block slides down a frictionless incline at $\theta=30^{\circ}$ while the
\emph{incline itself} is free to slide on a frictionless floor. Qualitatively:
does the block's acceleration exceed or fall short of $g\sin30^{\circ}$?
Why? (F=ma)
\end{exercise}

\section{Circular dynamics: the net force bends}\label{sec:new-circ}

Section 1.5 showed circular motion \emph{needs} $a_c=v^{2}/r$ inward;
Newton's second law says the \emph{net} force must supply it. The mantra:
$\sum F_{\mathrm{inward}}=mv^{2}/r$ --- centripetal is what the net force
\emph{does}, never an entry on the FBD.

\begin{example}[Banked curve, two boundary cases]\label{ex:banked}
A curve of radius $r=80\,\mathrm{m}$ is banked at
$\theta=15^{\circ}$. (a) What speed rounds it with \emph{no} friction?
(b) With $\mu_s=0.4$, what are the max and min speeds?
\par\smallskip
\textbf{Strategy.} (a) FBD has only $mg$ and $N$; resolve into vertical
(no accel) and horizontal (centripetal). (b) Friction acts along the
surface --- down-slope at high speed (fighting the tendency to slide
\emph{up} the bank), up-slope at low speed.
\begin{solution}
(a) Vertical: $N\cos\theta=mg$. Horizontal: $N\sin\theta=mv^{2}/r$.
Divide:
\[
  v=\sqrt{rg\tan\theta}=\sqrt{80\times9.8\times0.268}
   =14.5\,\mathrm{m/s}\approx52\,\mathrm{km/h}.
\]
(b) At high speed, friction points down-slope and adds an inward
component; the max-speed condition becomes (Morin's derivation)
$v_{\max}^{2}=rg\,\frac{\sin\theta+\mu_s\cos\theta}{\cos\theta-\mu_s\sin\theta}
=80\times9.8\times\frac{0.259+0.4\times0.966}{0.966-0.4\times0.259}$,
giving $v_{\max}=\sqrt{784\times0.747}=24.2\,\mathrm{m/s}\approx87$\,km/h.
The design window $14.5$--$24.2$\,m/s is what the bank buys. \checkmark
\end{solution}
\end{example}

\begin{figure}[htbp]\centering
\begin{tikzpicture}[scale=1.1,force/.style={->,thick,accent}]
\draw[thick] (0,0) -- (4.6,0) -- (4.6,2.0) -- cycle;
\draw (4.0,0) arc (180:180-24:0.6);
\node[font=\scriptsize] at (3.65,0.2) {$\theta$};
\filldraw[fill=gray!30] (2.0,0.92) rectangle (2.7,1.35);
\node[font=\scriptsize] at (2.35,1.14) {car};
\draw[force] (2.35,1.14) -- (2.35,-0.2) node[below,font=\scriptsize]{$m\mathbf g$};
\draw[force] (2.7,1.3) -- (3.35,2.0) node[above right=-2pt,font=\scriptsize]{$\mathbf N$};
\draw[force,amber] (2.0,0.95) -- (1.2,0.6) node[below left=-2pt,font=\scriptsize]{$\mathbf f$};
\draw[->,dashed] (5.0,1.2) -- (6.2,1.2);
\begin{scope}[xshift=5.0cm,yshift=0.6cm]
  \node[font=\scriptsize,align=left] at (1.4,0.9) {centre of circle\\lies to the left};
  \draw[->,thick,amber] (0.2,0.5) -- (-0.7,0.5) node[left,font=\scriptsize]{$a_c$};
\end{scope}
\end{tikzpicture}

\caption{Banked curve. Notice the normal force leans inward --- the bank's
job is to lend a horizontal component to $N$ so friction has less to do.}
\label{fig:banked}
\end{figure}

% from-zero:naive-net:begin
\begin{example}[First attempt: using the push instead of the net force]\label{ex:naive-net}
A $10\,\mathrm{kg}$ crate is pushed across a floor with a horizontal force of $30\,\mathrm{N}$; sliding friction
is $20\,\mathrm{N}$. A tempting calculation: $a=F/m=30/10=3\,\mathrm{m/s^2}$.
\par\smallskip
\textbf{What goes wrong.} Newton's second law uses the \emph{net} force, the vector sum of every force on the crate.
Friction opposes the motion, so the net force along the floor is $30-20=10\,\mathrm{N}$ and
$a=10/10=1\,\mathrm{m/s^2}$. The tempting answer was three times too large. This is why the free-body diagram
comes first: list every force, then add.
\end{example}
% from-zero:naive-net:end

\section{Friction: the force that adapts}\label{sec:new-friction}

Static friction \emph{matches} applied force up to a ceiling:
$f_s\le\mu_sN$; kinetic friction is a fixed drag $f_k=\mu_kN$ opposing
sliding, with $\mu_k<\mu_s$. Figure~\ref{fig:friction} shows the
characteristic: a 45-degree ramp that saturates, then drops to the kinetic
plateau. The \emph{angle of repose}\index{angle of repose} --- the steepest
angle before slipping --- sets $\tan\theta_{\max}=\mu_s$, a fast way to
measure $\mu_s$ with a protractor and patience.

\begin{figure}[htbp]\centering
\begin{tikzpicture}
\begin{axis}[width=8.5cm,height=5.5cm,xlabel={applied force $F$ (N)},
ylabel={friction force $f$ (N)},xmin=0,xmax=60,ymin=0,ymax=35,
xtick={0,20,40,60},ytick={0,10,20,30},legend pos=north west,font=\small]
\addplot[thick,accent,domain=0:30] {x};
\addplot[thick,accent,domain=30:60] {20};
\addplot[dashed,amber] coordinates {(30,0) (30,30)};
\addplot[mark=*,amber] coordinates {(30,30)} node[above right,font=\scriptsize]{$\mu_s N$};
\node[font=\scriptsize] at (axis cs:42,23) {$\mu_k N$ (kinetic)};
\node[font=\scriptsize,rotate=36] at (axis cs:12,15) {$f=F$ (static match)};
\end{axis}
\end{tikzpicture}

\caption{Static friction matches, saturates at $\mu_sN$, then drops to the
kinetic level $\mu_kN$. Notice the discontinuity: ``friction is weaker once
you're moving'' is why antilock brakes exist.}
\label{fig:friction}
\end{figure}

\begin{checkpoint}[Friction direction is derived, not guessed]
Two stacked blocks, the bottom one pushed: which way does friction act on
the \emph{top} block? If you said ``opposes motion'', you've guessed ---
friction on the top block points \emph{forward}: it's the only horizontal
force on it, and it's what accelerates it.
\end{checkpoint}

\begin{example}[Conical pendulum]\label{ex:conical}
A $0.5\,\mathrm{kg}$ mass on a $1.2\,\mathrm{m}$ string swings in a
horizontal circle at $30^{\circ}$ from vertical. Find its speed and period.
\par\smallskip
\textbf{Strategy.} FBD has $mg$ and $T$ only; $T$'s horizontal component is
the centripetal supply. Radius is $L\sin\theta$, not $L$.
\begin{solution}
Vertical: $T\cos\theta=mg\Rightarrow T=5.66\,\mathrm{N}$.
Horizontal: $T\sin\theta=mv^{2}/r$ with $r=1.2\sin30^{\circ}=0.6$\,m:
$v=\sqrt{Tr\sin\theta/m}=\sqrt{5.66\times0.6\times0.5/0.5}=1.84\,\mathrm{m/s}$.
Period $\tau=2\pi r/v=2.05\,\mathrm{s}$ --- independent of the mass, a
Galileo-clock property. \checkmark
\end{solution}
\end{example}

\section{Drag: terminal velocity as a fixed point}\label{sec:new-drag}

Drag opposes velocity with a strength that grows with speed: $cv$ (linear,
small spheres at low Reynolds number) or $cv^{2}$ (quadratic, everything
macroscopic in air or water). With quadratic drag,
\[
  m\dot v=mg-c_dv^{2}
\]
is an ODE with a \emph{fixed point}: at $v_t=\sqrt{mg/c_d}$ the right side
vanishes and the fall settles. The closed form
\[
  v(t)=v_t\tanh\!\left(\frac{gt}{v_t}\right)
\]
reaches $76\%$ of $v_t$ at $t=v_t/g$ --- for an 80-kg skydiver with
$c_d=0.30$, $v_t=51.1\,\mathrm{m/s}$, the answer a real dive converges to.

\begin{example}[Deriving $v(t)$ for quadratic drag]\label{ex:dragderiv}
Separate variables in $m\dot v=mg-c_dv^{2}$ and integrate.
\par\smallskip
\textbf{Strategy.} Write it as
$\frac{dv}{g-cv^{2}/m}=dt$ and substitute $v_t^{2}=mg/c_d$ so the
denominator is $g(1-v^{2}/v_t^{2})$ --- a standard form giving
$\operatorname{artanh}$.
\begin{solution}
\[
  \int\frac{dv}{1-v^{2}/v_t^{2}}=g\int dt
  \;\Rightarrow\;
  v_t\operatorname{artanh}\frac{v}{v_t}=gt
  \;\Rightarrow\;
  v(t)=v_t\tanh\frac{gt}{v_t}. \quad\checkmark
\]
The listing below confirms the derivation numerically --- Euler-Cromer
converges to the analytic curve.
\end{solution}
\end{example}

\begin{lstlisting}[language=Python,caption={\texttt{code/ch02\_drag\_fall.py} --- Euler-Cromer on quadratic drag vs the tanh solution.},label={lst:drag}]
def eulercromer(cd_over_m, dt, t_end):
    v, x, t = 0.0, 0.0, 0.0
    while t < t_end:
        v += (G - cd_over_m * v * v) * dt   # update velocity first
        x += v * dt                          # then position with NEW v
        t += dt
    return v, x
\end{lstlisting}

Output of a real run ($m=80$\,kg, $c_d=0.30$\,kg/m):
\begin{lstlisting}[numbers=none]
terminal velocity analytic: 51.1208 m/s
euler-cromer dt=0.5   v(20s)=51.0967 m/s  err=+0.0237  fall=858.0 m
euler-cromer dt=0.1   v(20s)=51.0785 m/s  err=+0.0055  fall=841.8 m
euler-cromer dt=0.01  v(20s)=51.0736 m/s  err=+0.0006  fall=838.1 m
at t=vt/g=5.22s: v=38.9333 m/s = 76.2% of vt
v=  5 m/s: linear drag a=1.000 m/s^2, quadratic a= 0.094 m/s^2
v= 20 m/s: linear drag a=4.000 m/s^2, quadratic a= 1.500 m/s^2
v= 50 m/s: linear drag a=10.000 m/s^2, quadratic a= 9.375 m/s^2
\end{lstlisting}

Notice the linear/quadratic crossover in the last lines: linear drag
dominates slow objects, quadratic dominates fast ones --- a skydiver's
$v_t$ sits in the quadratic regime (figure~\ref{fig:drag}).

\begin{figure}[htbp]\centering
\begin{tikzpicture}
\begin{axis}[width=8.5cm,height=5.5cm,xlabel={$v$ (m/s)},
ylabel={drag deceleration $a_{\rm drag}$ (m/s$^2$)},xmin=0,xmax=55,ymin=0,ymax=12,
xtick={0,10,20,30,40,50},legend pos=north west,font=\small]
\addplot[thick,accent,domain=0:55] {0.00375*x^2};
\addlegendentry{quadratic $cv^2/m$}
\addplot[thick,amber,domain=0:55] {0.2*x};
\addlegendentry{linear $cv/m$}
\addplot[dashed] coordinates {(51.12,0) (51.12,9.8)};
\addplot[dashed] coordinates {(0,9.8) (51.12,9.8)};
\addplot[mark=*] coordinates {(51.12,9.8)} node[below left,font=\scriptsize]{$v_t$: drag$=g$};
\end{axis}
\end{tikzpicture}

\caption{Drag deceleration vs speed for the two models. Terminal velocity
is the fixed point where drag equals $g$ (marked, $v_t=51.1$\,m/s).}
\label{fig:drag}
\end{figure}

\section{In engineering practice: load paths}\label{sec:new-eng}

Every load-bearing design --- bridge truss, building frame, airframe ---
is a walk through this chapter's FBD chain: gravity enters at the top,
travels through members as tension and compression, and exits into the
ground.
Structural engineers literally draw the discipline of
section~\ref{sec:new-fbd} as \emph{free-body cuts}: isolate a joint, draw
the member forces, solve $\sum F=0$, repeat. The banked-curve analysis is
the same trick highway designers run for every off-ramp, and the
friction ceiling of figure~\ref{fig:friction} is why ABS brakes pump:
holding the tire near $\mu_s$, not sliding on the weaker $\mu_k$ plateau.

\section{Common mistakes}\label{sec:new-mistakes}

\begin{tripwire}[$mv^{2}/r$ drawn as a force]
It is the \emph{required net inward force}. If your FBD shows $mg$, $N$,
\emph{and} a centripetal arrow, double-count.
\end{tripwire}
\begin{tripwire}[$N=mg$ on accelerating surfaces]
$N$ is determined by the constraint, not by gravity. Elevator
(example~\ref{ex:apparent}) and inclines (figure~\ref{fig:incline}) both
violate $N=mg$.
\end{tripwire}
\begin{tripwire}[Guessing the friction direction]
Friction opposes \emph{impending sliding at the contact} --- which can be
forward, as the checkpoint's stacked blocks show. Derive it from the
constraint, never from intuition about ``motion''.
\end{tripwire}
\begin{tripwire}[Third-law pairs on one body]
The pair acts on \emph{different} bodies. ``Normal cancels gravity'' is not
third-law bookkeeping --- those two are both on the block.
\end{tripwire}
\begin{tripwire}[Components before the FBD]
Picking axes and decomposing first is backwards: enumerate the real forces
on the real body, \emph{then} choose axes that simplify the algebra.
\end{tripwire}

\section{Summary}\label{sec:new-summary}

\begin{table}[ht]\centering
\begin{tabularx}{\linewidth}{@{}lY@{}}
\toprule
\tablehead\thead{Idea} & \thead{One-line version}\\
\midrule
three laws & inertial frames; $\sum\mathbf F=m\mathbf a$; pairs on different bodies \\
FBD method & isolate, enumerate contacts, pick axes, write $\sum F$ per axis \\
inclines & $mg\sin\theta$ drives, $N=mg\cos\theta$ only when surface is still \\
constraints & rope length $\Rightarrow$ equal-$a$; contact $\Rightarrow$ $N\ge0$ \\
circular dynamics & $\sum F_{\rm inward}=mv^{2}/r$; centripetal is a role, not a force \\
friction & $f_s$ matches up to $\mu_sN$; $f_k=\mu_kN$; $\tan\theta_{\rm repose}=\mu_s$ \\
drag & $cv$ or $cv^{2}$; fixed point $v_t=\sqrt{mg/c_d}$, $v=v_t\tanh(gt/v_t)$ \\
\bottomrule
\end{tabularx}
\caption{Newton's laws cheat sheet.}
\label{tab:newsummary}
\end{table}

\begin{wins}[Definition of done]
You can now: draw a correct FBD for any of the five canonical cases; solve
Atwood-style constraint problems; get the apparent weight in accelerating
frames; apply $\sum F_{\rm inward}=mv^{2}/r$ to curves and pendulums;
predict friction direction from the constraint; and derive + numerically
verify terminal velocity under quadratic drag.
\end{wins}

\chapternotes{%
The three-laws presentation and the FBD discipline follow MIT 8.01SC weeks
2--4 \cite{ocw801} and AP Physics C unit 2 \cite{apced}. OpenStax College
Physics chapters 4--6 \cite{openstax1} cover the same material with a large
worked set. The banked-curve and conical-pendulum examples follow Morin
chapters 1--3 \cite{morin}; the friction characteristic and repose angle are
standard \cite{hrw}. Drag models and the tanh solution follow 8.01SC's
resistive-media module; the skydiver numbers are representative, not
measured.}

\begin{rulebox}[{Watch instead}]
\begin{enumerate}
\item \emph{Newton's Laws} --- Walter Lewin, MIT 8.01x Lecture 6, 49:15,
\url{https://youtu.be/oduZsA0Tk58} --- \S\ref{sec:new-laws}--\S\ref{sec:new-circ}
at full length, same notation (accessed 2026-10-11).
\item \emph{Types of forces and free body diagrams} --- Khan Academy, 8:03,
\url{https://youtu.be/MaabUHLIIXA} --- the §\ref{sec:new-fbd} discipline in
eight minutes (accessed 2026-10-11).
\end{enumerate}
\end{rulebox}

\section*{Further reading}
\addcontentsline{toc}{section}{Further reading}

\begin{itemize}
\item MIT OpenCourseWare, \emph{8.01SC Classical Mechanics}, weeks 2--4 ---
Newton's laws, FBDs, friction and drag with recitation solutions. Free at
\url{https://ocw.mit.edu}, course ``8.01SC'' (verified 2026-10-11).
\item OpenStax, \emph{College Physics 2e}, ch.~4 ``Dynamics'' --- free text,
a second pass at the same laws with more worked FBDs. At
\url{https://openstax.org}, search ``College Physics 2e'' (verified 2026-10-11).
\item College Board, \emph{AP Physics C: Mechanics CED}, Unit 2 --- the exam
blueprint for this chapter's FR style. At
\url{https://apstudents.collegeboard.org} (verified 2026-10-11).
\item D.~Morin, \emph{Introduction to Classical Mechanics}, chs.~1--3 ---
FBDs, statics and accelerated frames at the level this book is aiming for;
ISBN 978-0521876223.
\item The Physics Classroom, \emph{Newton's Laws} tutorial --- readable
remediation for \S\ref{sec:new-laws}--\S\ref{sec:new-friction}, at
\url{https://www.physicsclassroom.com/class/newtlaws} (verified 2026-10-11).
\item Khan Academy, \emph{Forces and Newton's laws of motion} --- drills
matching every section here, at
\url{https://www.khanacademy.org/science/physics} (verified 2026-10-11).
\end{itemize}

\section*{Additional exercises}
\addcontentsline{toc}{section}{Additional exercises}

\begin{exercise}[easy]
Which of these belongs on a block-on-table FBD: (a) gravity on the block,
(b) the block's pull on the table, (c) friction on the block, (d) the
block's weight on the table. Justify each keep/drop. (AP-C)
\end{exercise}
\begin{exercise}[easy]
When does $N=mg\cos\theta$ hold on an incline? Give a case where it fails.
(AP-C)
\end{exercise}
\begin{exercise}[easy]
A $70\,\mathrm{kg}$ engineer weighs ``70\,kg'' on Earth. Convert to SI force;
then compute what a scale reads on the Moon ($g=1.6\,\mathrm{m/s^{2}}$)
and inside an elevator accelerating down at $3\,\mathrm{m/s^{2}}$. (AP-C)
\end{exercise}
\begin{exercise}[easy]
A rope holds a $10\,\mathrm{kg}$ mass static. Find the tension; find it
again while the mass accelerates up at $2\,\mathrm{m/s^{2}}$. (F=ma)
\end{exercise}
\begin{exercise}[easy]
A car's brakes can produce deceleration $0.8g$ on dry pavement
($\mu_k=0.8$). Estimate $\mu_k$ on ice if the same car needs $5\times$ the
dry distance to stop. (AP-C)
\end{exercise}
\begin{exercise}[medium]
Two stacked blocks: $m=2\,\mathrm{kg}$ on $M=5\,\mathrm{kg}$, frictionless
floor, $\mu_s=0.4$ between them. What is the largest horizontal force on
$M$ for which $m$ does not slip? (F=ma)
\end{exercise}
\begin{exercise}[medium]
Conical pendulum redux: keep $L$ and $\theta$ of
example~\ref{ex:conical} but double the mass. Which of $v$, $T$, $\tau$
change, and how? (AP-C)
\end{exercise}
\begin{exercise}[medium]
Atwood on a table: $m_1=2\,\mathrm{kg}$ on a frictionless table linked over
a pulley to hanging $m_2=3\,\mathrm{kg}$. Find $a$ and $T$; repeat with
$\mu_k=0.2$ on the table. (AP-C)
\end{exercise}
\begin{exercise}[medium]
The elevator of example~\ref{ex:apparent} accelerates at
$+2\,\mathrm{m/s^{2}}$ for $3$\,s, cruises, then brakes at
$-3\,\mathrm{m/s^{2}}$ to rest. Sketch the scale reading vs.\ time, and
compute the peak reading. (F=ma)
\end{exercise}
\begin{exercise}[medium]
Show that on the banked curve of example~\ref{ex:banked}, the minimum speed
with friction is
$v_{\min}^{2}=rg\,\frac{\sin\theta-\mu_s\cos\theta}{\cos\theta+\mu_s\sin\theta}$,
and evaluate it. What happens when $\mu_s\ge\tan\theta$? (F=ma)
\end{exercise}
\begin{exercise}[medium]
A skydiver with $c_d/m=0.004\,\mathrm{m^{-1}}$ reaches $v_t$; estimate
$v_t$, the time to $90\%\,v_t$, and the distance fallen by then. (APhO-style)
\end{exercise}
\begin{exercise}[hard]
A rope of mass $m$ and length $L$ lies on a frictionless table with a
length $x$ hanging over the edge. Show its acceleration is $a=gx/L$ and
find $x(t)$ --- the classic variable-length element problem. (USAPhO-style)
\end{exercise}
\begin{exercise}[hard]
A wedge of mass $M$ and angle $\theta$ is free to slide on a frictionless
floor; a block $m$ slides on its frictionless face. Set up the four
equations (two bodies, two axes + constraint) and solve for both
accelerations. (APhO-style)
\end{exercise}
\begin{exercise}[hard]
Estimate the maximum range of a projectile under quadratic drag:
linearise for small $c_d v_0^{2}/(mg)$ and show the leading correction to
the no-drag range is $-\tfrac{4}{3}c_d v_0^{3}/(m^{2}g^{2})\cdot(\text{shape
factor})$. Compare with the 36.6\% loss the ch01 code printed at
$k=0.02$. (USAPhO-style)
\end{exercise}
\begin{exercise}[medium]
Rederive the angle-of-repose rule: show that a block on an incline starts
sliding exactly when $\tan\theta=\mu_s$, and use it to compute $\mu_s$ if
sliding starts at $22^{\circ}$. (AP-C)
\end{exercise}
